\section{Egenværdier}
\subsection{Definition}
Givet vektorrum $V$ på $\mathbb{F}$ og lineær afbildning $L: V \to V$ 
er en skalar $\lambda \in \mathbb{F}$ en \textbf{egenværdi} for $L$, hvis der findes en vektor $\vec{v} \in V$, $\vec{v} \neq 0$, sådan at
\[L(v) = \lambda \cdot \vec{v}.\]
Vektoren $\vec{v}$ er en \textbf{egenvektor} tilhørende egenværdien $\lambda$.

\bemærk{$\lambda$ må godt være $0$}
\kilde{Definition 12.1.1}

\subsubsection{Matricers egenværdier}
Givet matrix $A \in \mathbb{F}^{n \times n}$ og vektor $\vec{v} \in \{ \mathbb{F}^n \setminus 0 \}$, er skalar $\lambda \in \mathbb{F}$ en \textbf{egenværdi} for $A$, hvis
\[A \vec{v} = \lambda \cdot \vec{v}.\]
Vektoren $\vec{v}$ er en \textbf{egenvektor} tilhørende egenværdien $\lambda$.

\bemærk{Gælder kun for kvadratiske matricer}
\kilde{Definition 12.1.2}

\subsection{Egenværdiproblemet}
For matrix $A \in \mathbb{F}^{n \times n}$, vil skalar $\lambda \in \mathbb{F}$ være en egenværdi for $A \iff$
\[\det(A - \lambda \cdot \mathbf{I}_n) = 0.\]
Dette er det \textbf{karakteristiske polynomium} for $A$ og betegnes $p_{\mathbf{A}}(\lambda)$.

\

Med andre ord, rødderne i følgende polynomium er egenværdierne for $\mathbf{A}$.
\[
\det \left(
    \begin{bmatrix}
        a_{11} - \lambda & a_{12} & \cdots & a_{1n} \\
        a_{21} & a_{22} - \lambda & \cdots & a_{2n} \\
        \vdots & \vdots & \ddots & \vdots \\
        a_{n1} & a_{n2} & \cdots & a_{nn} - \lambda 
    \end{bmatrix}\right) = 0
\]

Dette vil give et polynomium af grad $n$, som løses for det ubekendte $\lambda$.

\kilde{Sætning 12.1.1}

\subsubsection{Egenværdi for Lineære afbildinger}
For en lineær afbildning $L: V \to V$ og en basis $\beta$ for $V$, vil en skalar $\lambda \in \mathbb{F}$ være en egenværdi for $L$ hvis og kun hvis
\[
\det(_\beta[L]_{\beta} - \lambda \cdot \mathbf{I}_n) = 0.
\kildetag{Sætning 12.1.2}
\]

\paragraph{Om egenværdier af lineære afbildninger}

\

Den valgte basis er ligegyldig.

\begin{align*}
    p_{_\beta[L]_\beta}(\lambda) &= p_{_\gamma[L]_\gamma}(\lambda) \kildetag{Sætning 12.1.4} \\
    \det(_\beta[L]_\beta) &= \det(_\gamma[L]_\gamma) \kildetag{Korollar 12.1.5}
\end{align*}

Det karakteristiske polynomium for en lineær afbildning $L: V \to V$ betegnes med valgfri basis:
\[
p_L(\lambda) = \det(_\beta[L]_\beta - \lambda \cdot \mathbf{I}_n) \kildetag{Definition 12.1.3}
\]

\subsection{Egenrum}

Egenrummet for en egenværdi $\lambda$ og matrix $\mathbf{A} \in \mathbb{F}^{n \times n}$ betegnes $E_{\lambda}$ og defineres som mængden af alle egenvektorer tilhørende $\lambda$ samt \emph{nulvektoren}. Det kan findes ved at finde kernen af matricen:

\[ 
\begin{bmatrix}
    a_11 - \lambda & a_{12} & \cdots & a_{1n} \\
    a_{21} & a_{22} - \lambda & \cdots & a_{2n} \\
    \vdots & \vdots & \ddots & \vdots \\
    a_{n1} & a_{n2} & \cdots & a_{nn} - \lambda
\end{bmatrix} \kildetag{Lemma 12.2.3}
\]

Se evt. sektion \ref{sec:kerne} om hvordan kernen findes.

\subsubsection{Egenrum er underrum}
\begin{align*}
    E_\lambda = \{ \vec{v} \in V | L(\vec{v}) = \lambda \cdot \vec{v} \} \kildetag{Sætning 12.2.1} \\
    A_\lambda = \{ \vec{v} \in V |A \cdot \vec{v} = \lambda \cdot \vec{v} \} \kildetag{Sætning 12.2.1}
\end{align*}
Disse er underrum i hhv. $V$ og $A$.

\section{Diagonalisering}

\subsection{Kan en matrix diagonaliseres?}
Lineær afbildning $L: V \to V$ eller matrix $\mathbf{A} \in \mathbb{F}^{n \times n}$ er diagonaliserbar $\iff$ der findes en basis for $V$ bestående af egenvektorer tilhørende $L$.

Det skal gælde for alle egenværdier at $\gm(\lambda) = \am(\lambda)$.

Det er det samme som at det karakteristiske polynomium for $L$ kan skrives som
\[
p_L(Z) = (-1)^n \cdot (Z - \lambda_1)^{m_1} \cdots (Z - \lambda_r)^{m_r}
\kildetag{Sætning 12.3.3}
\]

Når en afbildning kan diagonaliseres kan man finde en diagonalmatrix $\mathbf{D}$ og en invertibel matrix $\mathbf{Q}$ så
\[
L = \mathbf{Q} \cdot \mathbf{D} \cdot \mathbf{Q}^{-1} \kildetag{Lemma 12.4.2}
\]
Se sektion \ref{sec:diagonal_matrix} og \ref{sec:invertibel_matrix}.

\subsection{Find Diagonalmatrix}
\label{sec:diagonal_matrix}
Hvis $\mathbf{A} \in \mathbb{F}^{n \times n}$ kan diagonaliseres, kan diagonalmatricen findes ved at indsætte egenværdierne for $\mathbf{A}$ på diagonalen i en $n \times n$ matrix $\mathbf{D}$.
\[\mathbf{D} = 
\begin{bmatrix}
    \lambda_1 & 0 & \cdots & 0 \\
    0 & \lambda_2 & \cdots & 0 \\
    \vdots & \vdots & \ddots & \vdots \\
    0 & 0 & \cdots & \lambda_n
\end{bmatrix}\]

\bemærkBig{Hvis en egenværdi har algebraisk multiplicitet $n$, skal den optræde $n$ gange efter hinanden på diagonalen i $\mathbf{D}$.}

\subsection{Find invertibel matrix}
\label{sec:invertibel_matrix}
Diagonalmatrix $\mathbf{D}$ kan skrives som $\mathbf{D} = \mathbf{Q}^{-1} \cdot \mathbf{A} \cdot \mathbf{Q}$, hvor $\mathbf{Q}$ er en invertibel matrix bestående af egenvektorerne for egenværdierne i $\mathbf{D}$ som søjler.
\[\mathbf{Q} =
\begin{bmatrix}
    | & | & & | \\
    \vec{v_1} & \vec{v_2} & \cdots & \vec{v_n} \\
    | & | & & |
\end{bmatrix}\]

\bemærkBig{Rækkefølgen af egenvektorerne i $\mathbf{Q}$ skal svare til rækkefølgen af egenværdierne i $\mathbf{D}$.}


\subsubsection{Algebraisk multiplicitet}
$\am(\lambda)$ er antallet af gange egenværdi $\lambda$ optræder i roden af $p_\mathbf{A}(\lambda)$.

\eksempel{
$\lambda^{10} \cdot (\lambda-5)^6 \cdot (\lambda+8)^{10} = 0$
$ \implies \am(-8) = 10, \am(0) = 10, \am(5) = 6$.
}

\subsubsection{Geometrisk multiplicitet}
\[\gm(\lambda) = \dim(E_\lambda)\]

\eksempel{
Givet
\(
E_{-3} = \spanop \left\{
s \cdot 
\begin{bmatrix}
1 \\ 0 \\ 0
\end{bmatrix} + t \cdot 
\begin{bmatrix}
    0 \\ 1 \\ 0
\end{bmatrix}
\right\}
\)
er $\gm(-3) = 2$.
}

\subsubsection{Generel sammenhæng}
\[
1 \le \gm(\lambda) \le \am(\lambda) \le \dim(V)
\kildetag{Sætning 12.2.4}
\]

\subsection{Koordinater mht. egenvektor, fra egenværdi}
Hvis der gives to vektorer $\vec{v_1}, \vec{v_2} \in V$ og en lineær afbildning $f: V \to V$, og det angives at $\vec{v_1}$ er en egenvektor tilhørende egenværdien $\lambda_1$ og $\vec{v_2}$ er en egenvektor tilhørende egenværdien $\lambda_2$, vil koordinaterne for \(f(2 \cdot \vec{v_1} - 4 \cdot \vec{v_2})\) med hensyn til basis \(\beta = \{\vec{v_1}, \vec{v_2}\}\) være:
\[
    (2 \cdot \lambda_1, \quad -4 \cdot \lambda_2)
\]

\subsection{Find Diagonalmatrix fra egenvektor matrix}

Gives en matrix \( \mathbf{A} = \begin{bmatrix}
    a & b \\
    c & d
\end{bmatrix}
\) og en invertibel matrix \( \mathbf{Q} = \begin{bmatrix}
    w & x \\
    y & z
\end{bmatrix}
\) bestående af egenvektorerne til $\mathbf{A}$ som søjler, kan diagonalmatrix $\mathbf{D}$ fines ved at gange hver af søjlerne i $\mathbf{Q}$ med $\mathbf{A}$, og se hvilket multiplum af søjlen man får.


\subsection{Andet}

For $\mathbf{A}^{2 \times 2}$ med kendt determinant og en egenværdi er den anden egenværdi givet ved
\[\lambda_2 = \frac{\det(\mathbf{A})}{\lambda_1}.\]

\

$x$ forskellige egenværdier $\implies$ mindst $x$ lineært uafhængige egenvektorer.

\

Egneværdierne for $\mathbf{A}^{-1}$ er givet ved $\frac{1}{\lambda_i}$ hvor $\lambda_i$ er egenværdierne for $\mathbf{A}$.

\

2 lineært uafhængige egenvektorrum har kun én fælles vektor, nemlig nulvektoren.